% CS 418: Intro to Data Science — Worksheet 15 (15-week08-day1)
% Week 8, Day 1
% NOTICE: This worksheet was drafted by an LLM coding system from the
% instructor's choice of checkpoints. It is maintained, reviewed, and vetted by humans.
% Compile with: latexmk -pdf worksheet.tex

\documentclass[11pt]{article}

% ── Packages ──────────────────────────────────────────────────────────────────
\usepackage[margin=0.65in, headheight=14pt]{geometry}
\usepackage{amsmath, amssymb}
\usepackage{ragged2e}
% Never break a word across lines, including at hyphens ("Sans-Serif").
\hyphenpenalty=10000
\exhyphenpenalty=10000
\usepackage{enumitem}
\usepackage{booktabs}
\usepackage{graphicx}
\usepackage{hyperref}
\usepackage{xcolor}
\usepackage{fancyhdr}
\usepackage{mdframed}
\usepackage{array}

% ── Page style ────────────────────────────────────────────────────────────────
\pagestyle{fancy}
\fancyhf{}
\lhead{\textbf{CS 418: Intro to Data Science}}
\rhead{UIC, Fall 2026}
\cfoot{\thepage}
\renewcommand{\headrulewidth}{1pt}
\setlength{\headheight}{12pt}
\setlength{\headsep}{0.4cm}

% ── Hyperlinks ────────────────────────────────────────────────────────────────
\hypersetup{colorlinks=false, hidelinks}

% ── Convenience environments ──────────────────────────────────────────────────
% Usage: \blankspace{6cm} — blank area for hand-written work, no rules
\newcommand{\blankspace}[1]{%
  \vspace{0.3em}
  \begin{minipage}[t][#1][t]{\linewidth}
  \end{minipage}
  \vspace{0.3em}
}

% Usage: \guidedopen{title}{guided content}{guided workspace height}{open-ended workspace height}
% Left: a titled, structured prompt with its own blank workspace. Right: a
% fully blank workspace for open-ended exploration of the same topic (no
% prompt text), sized so the two columns land at roughly the same total
% height. Neither workspace is ruled.
\newcommand{\guidedopen}[4]{%
  \noindent\textbf{#1}\\[0.18em]
  \noindent
  \begin{minipage}[t]{0.485\textwidth}
    #2
    \blankspace{#3}
  \end{minipage}
  \hfill
  \begin{minipage}[t]{0.485\textwidth}
    \blankspace{#4}
  \end{minipage}
  \\[0.4em]
}

% Usage: \panelbox{width}{height}{caption} — a framed drawing panel with `caption`
% centered beneath it. An empty caption leaves the panel bare, with nothing
% below the frame at all.
\newcommand{\panelbox}[3]{%
  \begin{minipage}[t]{\dimexpr#1+2\fboxsep+2\fboxrule\relax}
    \framebox[#1][c]{\rule{0pt}{#2}}%
    \if\relax\detokenize{#3}\relax\else
      \\[0.25em]{\centering\small #3\par}%
    \fi
  \end{minipage}%
}

% Usage: \axespanel{width}{height}{caption} — a \panelbox with faint axes: the y-axis centered, the x-axis near the bottom.
\newcommand{\axespanel}[3]{%
  \begin{minipage}[t]{\dimexpr#1+2\fboxsep+2\fboxrule\relax}
    \framebox[#1][c]{\rule{0pt}{#2}%
      \makebox[0pt][c]{\textcolor{black!25}{\rule{0.6pt}{#2}}}%
      \makebox[0pt][c]{\textcolor{black!25}{\raisebox{\dimexpr#2/6\relax}{\rule{#1}{0.6pt}}}}}%
    \\[0.25em]{\centering\small #3\par}%
  \end{minipage}%
}

% ══════════════════════════════════════════════════════════════════════════════
\begin{document}
\RaggedRight

% ── Title block: topic left, info box right ───────────────────────────────────
\noindent
\begin{minipage}[t]{0.55\textwidth}
  \vspace{0pt}
  {\LARGE \textbf{Worksheet \#\,15}} \\[0.18em]
  {\large \textit{Week 8, Day 1: Linear regression}}
\end{minipage}
\hfill
\begin{minipage}[t]{0.40\textwidth}
  \vspace{0pt}
  \small
  Name: \underline{\hspace{4.5cm}} \\[0.35em]
  NetID: \underline{\hspace{4.3cm}} \\[0.35em]
  Date: \underline{\hspace{4.4cm}} \\
\end{minipage}

\vspace{0.6em}
\noindent\rule{\linewidth}{0.6pt}
\vspace{0.1em}

\noindent\underline{\textbf{Guided checkpoints}}\\[0.2em]
\noindent{\small For full credit, fill in the checkpoints. If you skip a checkpoint, you must include your own notes for full credit.}
\vspace{0.8em}

% ── 1. What is a model? ──────────────────────────────────────────────────────
\guidedopen{1. What Is a Model?}{
Write down a definition of a model, in your own words or from one of the examples on the slides.
}{1.2cm}{2.2cm}

\vfill

% ── 2. All models ────────────────────────────────────────────────────────────
\guidedopen{2. George Box}{
Write down George Box's quote about ``all models.''
}{0.8cm}{1.4cm}

\vfill

% ── 3. Mechanistic vs. statistical ───────────────────────────────────────────
\guidedopen{3. Mechanistic or Statistical?}{
Write down each example and whether it is mechanical or statistical.

\vspace{0.5em}
A.
\vspace{0.65cm}

B.
\vspace{0.65cm}

C.
\vspace{0.65cm}

D.
}{0.65cm}{5.4cm}

\vfill

% ── 4. Loss function ─────────────────────────────────────────────────────────
\guidedopen{4. Loss Functions}{
Write the average loss $L(\theta, \mathbf{y})$ from the slides, then explain it in plain language.

\vspace{0.5em}
Formula:
\vspace{1.2cm}

In plain language:
}{1.3cm}{4.4cm}

\vfill

\newpage

% ── 5. Squared vs. absolute error ────────────────────────────────────────────
\guidedopen{5. Squared vs. Absolute Error}{
Sketch each loss against the error $y - \theta$.

\vspace{0.45em}

\noindent
\axespanel{3.7cm}{2.7cm}{squared error}\hfill
\axespanel{3.7cm}{2.7cm}{absolute error}
}{0.2cm}{3.8cm}

\vfill

% ── 6. The modeling process ──────────────────────────────────────────────────
\guidedopen{6. The Modeling Process}{
List the main steps of modeling, in order.

\vspace{0.5em}
1.
\vspace{0.6cm}

2.
\vspace{0.6cm}

3.
\vspace{0.6cm}

4.
}{0.6cm}{4.0cm}

\vfill

% ── 7. Constant and simple linear ────────────────────────────────────────────
\guidedopen{7. Linear Models}{
Write the equation for each model.

\vspace{0.5em}
Constant model:
\vspace{0.8cm}

Simple linear regression:
}{0.8cm}{2.4cm}

\vfill

% ── 8. LINE ──────────────────────────────────────────────────────────────────
\guidedopen{8. LINE}{
Name each LINE criterion and what it requires.

\vspace{0.5em}
L:
\vspace{0.7cm}

I:
\vspace{0.7cm}

N:
\vspace{0.7cm}

E:
}{0.7cm}{4.2cm}

\vfill

% ── 9. Multiple linear ───────────────────────────────────────────────────────
\guidedopen{9. Multiple Linear Regression}{
Write the equation for multiple linear regression.
}{0.8cm}{1.4cm}

\vfill

\end{document}
